\documentclass{article}
\usepackage{iclr2024_conference,times}

\usepackage[utf8]{inputenc}
\usepackage[T1]{fontenc}
\usepackage{hyperref}
\usepackage{url}
\usepackage{booktabs}
\usepackage{amsfonts}
\usepackage{nicefrac}
\usepackage{microtype}
\usepackage{titletoc}

\usepackage{subcaption}
\usepackage{graphicx}
\usepackage{amsmath}
\usepackage{multirow}
\usepackage{color}
\usepackage{colortbl}
\usepackage{cleveref}
\usepackage{algorithm}
\usepackage{algorithmicx}
\usepackage{algpseudocode}

\DeclareMathOperator*{\argmin}{arg\,min}
\DeclareMathOperator*{\argmax}{arg\,max}

\graphicspath{{../}}
\begin{filecontents}{references.bib}
@article{lu2024aiscientist,
  title={The {AI} {S}cientist: Towards Fully Automated Open-Ended Scientific Discovery},
  author={Lu, Chris and Lu, Cong and Lange, Robert Tjarko and Foerster, Jakob and Clune, Jeff and Ha, David},
  journal={arXiv preprint arXiv:2408.06292},
  year={2024}
}

@article{nguyen2024evo2,
  title={Genome modeling and design across all domains of life with {Evo} 2},
  author={Nguyen, Eric and others},
  journal={Nature},
  year={2025}
}

@article{kanehisa2000kegg,
  title={{KEGG}: Kyoto encyclopedia of genes and genomes},
  author={Kanehisa, Minoru and Goto, Susumu},
  journal={Nucleic Acids Research},
  volume={28},
  number={1},
  pages={27--30},
  year={2000}
}

@article{cunningham2022ensembl,
  title={Ensembl 2022},
  author={Cunningham, Fiona and others},
  journal={Nucleic Acids Research},
  volume={50},
  number={D1},
  pages={D988--D995},
  year={2022}
}

@article{tsherniak2017defining,
  title={Defining a cancer dependency map},
  author={Tsherniak, Aviad and others},
  journal={Cell},
  volume={170},
  number={3},
  pages={564--576},
  year={2017}
}

@article{sondka2018cosmic,
  title={The {COSMIC} Cancer Gene Census: describing genetic dysfunction across all human cancers},
  author={Sondka, Zbyslaw and others},
  journal={Nature Reviews Cancer},
  volume={18},
  number={11},
  pages={696--705},
  year={2018}
}

@article{devlin2019bert,
  title={{BERT}: Pre-training of deep bidirectional transformers for language understanding},
  author={Devlin, Jacob and Chang, Ming-Wei and Lee, Kenton and Toutanova, Kristina},
  journal={arXiv preprint arXiv:1810.04805},
  year={2018}
}

@article{dalla2023nucleotide,
  title={The nucleotide transformer: building and evaluating robust foundation models for human genomics},
  author={Dalla-Torre, Hugo and others},
  journal={bioRxiv},
  year={2023}
}
\end{filecontents}

\title{Do Evo2 Embeddings Encode Gene Function? \\
Classification of DNA Repair vs.\ Tumor Suppressor Genes \\
from Promoter and Coding Sequences}

\begin{document}

\maketitle

\begin{abstract}
DNA foundation models trained on large genomic corpora have demonstrated strong performance on sequence-level tasks such as variant effect prediction and gene expression modeling. However, whether their internal representations implicitly encode higher-level biological function — without explicit functional supervision — remains an open question. We investigate this using Evo2, a 7-billion-parameter autoregressive DNA language model trained on 8.8 trillion nucleotide tokens. For a curated set of human genes from two functionally distinct classes — DNA repair genes (from KEGG pathways) and tumor suppressor genes (from the Cancer Gene Census) — we extract mean-pooled late-layer embeddings from the frozen model using each gene's 2kb proximal promoter concatenated with its canonical coding sequence. A shallow MLP classifier trained on these embeddings is used to test whether functional class is linearly separable in Evo2's representation space. We report classification accuracy, AUROC, and F1, and use this framework to systematically investigate which layers, sequence regions, and model components contribute most to functional discriminability.
\end{abstract}

\section{Introduction}
\label{sec:intro}

The central dogma of molecular biology connects DNA sequence to molecular function, yet the mapping is complex, non-local, and context-dependent. Promoter architecture, codon usage bias, regulatory element density, and evolutionary conservation all contribute to a gene's functional identity. Genomic foundation models such as Evo2 \citep{nguyen2024evo2} and the Nucleotide Transformer \citep{dalla2023nucleotide} are trained to predict masked or next tokens across billions of genomic sequences, and in doing so may learn to represent these patterns implicitly.

A natural question is whether these representations encode biologically interpretable structure beyond local sequence statistics. Specifically: can a model trained purely on next-nucleotide prediction distinguish between genes with fundamentally different cellular roles — such as maintaining genome integrity (DNA repair) versus suppressing oncogenic transformation (tumor suppression)?

These two classes make a compelling test case. DNA repair genes (e.g.\ \textit{MLH1}, \textit{BRCA1}, \textit{RAD51}) share conserved functional domains and co-regulatory elements. Tumor suppressor genes (e.g.\ \textit{TP53}, \textit{RB1}, \textit{PTEN}) are a more heterogeneous class united by loss-of-function oncogenic consequence rather than sequence similarity. If Evo2 embeddings separate these classes, it implies that the model has learned functionally relevant genomic grammar. If they do not, it points to limits of sequence-based representations for functional inference.

\section{Methods}
\label{sec:methods}

\subsection{Gene Sets}
DNA repair genes are retrieved from KEGG pathways hsa03430 (mismatch repair), hsa03420 (nucleotide excision repair), hsa03440 (homologous recombination), and hsa03410 (base excision repair). Tumor suppressor genes are drawn from a curated list derived from the COSMIC Cancer Gene Census \citep{sondka2018cosmic}.

\subsection{Sequence Retrieval}
For each gene, we retrieve the human genomic sequence via the Ensembl REST API \citep{cunningham2022ensembl}: specifically, 2000bp upstream of the transcription start site (proximal promoter) concatenated with the canonical coding sequence (longest annotated CDS). Sequences are uppercased and truncated to 8192 nucleotides to fit within Evo2's inference context.

\subsection{Evo2 Embeddings}
We load the frozen Evo2-7B model \citep{nguyen2024evo2} and extract embeddings from layer \texttt{blocks.28.mlp.l3} — a late-layer representation recommended by the original authors as capturing rich contextual information. Embeddings are mean-pooled over the sequence length dimension to produce a fixed-size gene-level representation.

\subsection{Classifier}
A two-layer MLP (hidden dimensions 256 and 128, ReLU activations, dropout 0.2) is trained with binary cross-entropy loss to classify DNA repair (label 0) vs.\ tumor suppressor (label 1). We use an 80/20 train/test split with three random seeds and report mean $\pm$ standard error.

\section{Results}
\label{sec:results}

\begin{figure}[h]
    \centering
    \begin{subfigure}{0.48\textwidth}
        \includegraphics[width=\textwidth]{train_loss.png}
        \caption{Training loss over epochs.}
    \end{subfigure}
    \hfill
    \begin{subfigure}{0.48\textwidth}
        \includegraphics[width=\textwidth]{val_auroc.png}
        \caption{Validation AUROC over epochs.}
    \end{subfigure}
    \caption{Learning curves for the Evo2 + MLP classifier, averaged over 3 seeds.}
    \label{fig:curves}
\end{figure}

\begin{figure}[h]
    \centering
    \includegraphics[width=0.9\textwidth]{final_metrics_bar.png}
    \caption{Final held-out classification metrics ($\pm$ stderr over 3 seeds).}
    \label{fig:final_metrics}
\end{figure}

\section{Conclusion}
\label{sec:conclusion}

We present a framework for probing whether Evo2 genomic embeddings encode functional gene class information, using binary classification of DNA repair vs.\ tumor suppressor genes from promoter and coding sequence. This template enables systematic investigation of which layers, sequence regions, and model scales contribute to functional discriminability — a key question for understanding what genomic foundation models learn during pretraining.

\bibliography{references}
\bibliographystyle{iclr2024_conference}

\end{document}
